\section{Renormalization of gluon-gluon-gauge-link vertex}
Finally, we exam the renormalization of UV divergence of the topology (d) in Fig.~\ref{fig:div-topo-g}, whose diagrams involve two gluons and a gauge link, and could be referred as the gluon-gluon-gauge-link vertex.  Similar to the Ward identity in Eq.~(\ref{eq:gwi1}), we construct the following Ward identity for the ``bare" fields and operators,
\begin{align}
&\langle \partial_\lambda^x A_d^\lambda(x) \partial_\rho^y A_e^\rho(y) [\Phi(\{\xi_{2z},\xi_{1z}\})]_{a b} \,  F_{b}^{\mu \nu}(\xi_{1z})  \rangle \nonumber \\
+&i \delta^{(d)}(x-y) \delta_{de} \langle [\Phi(\{\xi_{2z},\xi_{1z}\})]_{a b} \,  F_{b}^{\mu \nu}(\xi_{1z}) \rangle
\label{eq:gwi2} \\
=& g_{s} \langle f^{afg} \bar{c}_e (y) c_f (\xi_{2z}) \partial_\lambda^x A_d^\lambda(x) [\Phi(\{\xi_{2z},\xi_{1z}\})]_{g b} \,  F_{b}^{\mu \nu}(\xi_{1z}) \rangle.
\nonumber
\end{align}
A pictorial interpretation of Eq.~(\ref{eq:gwi2}) is given in Fig.~\ref{fig:gwi2}. Similar to Fig.~\ref{fig:gwi1}, the topology of the left-hand side of Fig.~\ref{fig:gwi2} is the same as that of the diagram (d) in Fig.~\ref{fig:div-topo-g}, except the external gluons of the diagrams are contracted with their respective momenta and expressed in coordinate space.  The right-hand side of the equation in Fig.~\ref{fig:gwi2} is UV finite after all previously discussed renormalizations performed, including QCD Lagrangian, gauge links, and gluon-gauge-link vertex. That is,  we find that the topology (d) in Fig.~\ref{fig:div-topo-g} is free of UV divergence if both external gluons are contracted by their respective momenta.

%%%%%%%%%%%%%%%%%%%%%%%%%%%%%%%
\begin{figure}[t]
	\begin{center}
		\includegraphics[width=0.55\textwidth]{images/Green-2-g.pdf}
		\caption{\label{fig:gwi2} Pictorial interpretation of the generalized Ward identity in Eq.~(\ref{eq:gwi2})
		with dashed line represents the ghost field.}
	\end{center}
\end{figure}
%%%%%%%%%%%%%%%%%%%%%%%%%%%%%%

Similar to the discussion of the gluon-gauge-link vertex, the Ward identity helps reduce the superficial UV divergence of the topology (d) in Fig.~\ref{fig:div-topo-g}. Since the diagrams of the topology (d) have only superficial logarithmic divergence, the additional reduction of the superficial UV divergence from the Ward identity makes the topology (d) in Fig.~\ref{fig:div-topo-g} UV finite. Therefore, after the renormalization of topology (c), the topology (d) requires no  additional renormalization.  This is similar to the case of renormalizing gluon vertexes of QCD Lagrangian, where gauge invariance guarantees that four-gluon vertex will be free of UV divergence once three-gluon vertex is renormalized.
